%% Согласно ГОСТ Р 7.0.11-2011:
%% 5.3.3 В заключении диссертации излагают итоги выполненного исследования, рекомендации, перспективы дальнейшей разработки темы.
%% 9.2.3 В заключении автореферата диссертации излагают итоги данного исследования, рекомендации и перспективы дальнейшей разработки темы.


В рамках данной работы был разработан пакет программ
для моделирования динамики жидкости численым методом SPH
с применением массового параллелизма для применения в суперкомпьютерах
разного типа (ориентированных на CPU или на GPU).
%
Исходный код проекта является открытым.
%
Ссылка на актуальную версию исходного кода на GitHub [Электронный ресурс]: 
\url{https://github.com/RackotRR/RRSPH} (дата обращения: 14.06.2025).

Рассмотрена и реализована математическая модель SPH для решения системы уравнений Навье-Стокса.
%
Рассмотрены подходы к моделированию несжимаемости в методе SPH,
реализован метод WCSPH (Weakly Compressible Smoothed Particle Hydrodynamics).
%
Реализован ряд сглаживающих функций,
различные методы моделирования плотности и сил давления.
%
Реализованы численные члены, такие как 
искусственная вязкость,
искусственное давление,
усреднение скорости XSPH.
%
Реализованы метод динамических частиц и метод отталкивающих частиц для обработки граничных условий.
%
Для интегрирования по времени применён метод Leapfrog с динамическим шагом, рассчитанным из условия устойчивости.
%
Реализовано несколько методов генерации волн с помощью динамических границ (поршня):
методы первого и второго порядка для моделирования волн,
метод Рэлэя для генерации солитонной волны.
%
Реализован сеточный метод поиска пар взаимодействующих частиц.

Выполнено распараллеливание с использованием технологий OpenMP и OpenCL для ускорения расчётов.
%
В разработке применена система построения CMake, 
что позволило собирать проект различными способами на различных системах --- как на ОС Windows, так и на ОС Linux.
%
Используется система контроля версий Git и модули git (часть библиотек, 
а также \verb|SPH2D_Drawer| включены в репозиторий как модули,
что позволяет разделить проект на несколько подпроектов с минимальным набором зависимостей).

Реализованы программы для анализа результатов моделирования, 
которые позволяют из данных в виде набора частиц получать информацию о значении функции в указанной точке (\verb|FuncAtPoint|), 
либо профиль жидкости, как по пространству, так и по времени (\verb|WaterProfile|).

Реализована программа для визуализации данных в виде частиц \verb|SPH2D_Drawer| в 2D.
%
Она обладает системой команд для настройки отображения в процессе визуализации, 
способна воспроизводить и записывать анимации. 
%
Такой инструмент, как тепловая карта, позволяет, помимо положения частиц, получать информацию о состоянии данной частицы 
(например, о давлении или о проекции скорости).
%
Реализована программа \verb|Scatter3D| для визуализации данных в 3D. 
%
Она обладает основными функциями \verb|SPH2D_Drawer|, 
но позволяет эффективно визуализировать миллионы частиц,
однако не поддерживает команды и тепловую карту.

Разработаны python-скрипты для подготовки начальных условий для моделирования
простых экспериментов с разрушением дамбы или генерацией волн. 
%
Разработан скрипт экспорта временных слоёв и параметров частиц из пакета \verb|DualSPHysics|,
что позволяет выполнять их обработку имеющимися средствами,
а также выполнять аналогичные расчёты.
%
Разрабатывается графический пользовательский интерфейс для удобной настройки параметров модели
и быстрого запуска рачётов.

Проведён анализ эффективности разработанного ПО.
%
Проведены численные эксперименты, 
позволяющие определить погрешности в дисперсионном соотношении созданных волн:
показана высокая точность модели.
%
Проведён анализ эффективной вязкости жидкости на основе показателя затухания волн в численной модели:
вязкость значительно превышает аналитическое значение, 
но близка к эффективной вязкости других моделей на основе SPH.
%
Результаты тестирования демонстрируют соответствие модели теоретическим ожиданиям и 
её применимость для широкого класса гидродинамических задач.

В ходе работы были освоены следующие компетенции:

УК-1 Способен осуществлять критический анализ проблемных ситуаций на основе системного подхода, вырабатывать стратегию действий.
Компетенция была сформирована в процессе:
\begin{enumerate}
    \item обсуждения работы с научным руководителем;
    \item определения требований к пакету программ для моделирования динамики жидкости;
    \item выбора библиотек и технологий для последующей разработки программного обеспечения на их основе;
    \item анализа алгоритмов и методов решения задач гидродинамики, описанных в научной литературе;
    \item поиска информации в различных академических источниках;
    \item оценки промежуточных результатов тестирования.
\end{enumerate}

УК-2 Способен управлять проектом на всех этапах его жизненного цикла.
Компетенция была сформирована в процессе:
\begin{enumerate}
    \item обсуждения работы с научным руководителем;
    \item определения требований к пакету программ для моделирования динамики жидкости;
    \item проектирования пакета программ для моделирования динамики жидкости, определения его архитектуры и модулей;
    \item реализации пакета программ для моделирования динамики жидкости;
    \item тестирования численных моделей, построенных с помощью разработанного ПО;
    \item составления отчётов и презентаций в рамках практик, выпускной квалификационной работы и конференций.
\end{enumerate}

УК-3 Способен организовывать и руководить работой команды, вырабатывая командную стратегию для достижения поставленной цели.
Компетенция была сформирована в процессе:
\begin{enumerate}
    \item обсуждения работы с научным руководителем;
    \item определения требований к пакету программ для моделирования динамики жидкости;
    \item исполнения требуемой роли на каждом этапе проекта;
    \item работа в команде при подготовке тезисов для конференций и текста выпускной квалификационной работы.
\end{enumerate}

УК-4 Способен применять современные коммуникативные технологии, в том числе на иностранном(ых) языке(ах), 
для академического и профессионального взаимодействия.
Компетенция была сформирована в процессе:
\begin{enumerate}
    \item подготовки тезисов для конференций;
    \item составления отчёта по выпускной квалификационной работе;
    \item защиты выпускной квалификационной работы;
    \item документирования исходного кода разрабатываемого ПО;
    \item изучения академической литературы на иностранном языке;
    \item изучения технической документации на иностранном языке;
    \item обсуждения работы с научным руководителем при личных встречах, либо с помощью средств телекоммуникации.
\end{enumerate}

УК-5 Способен анализировать и учитывать разнообразие культур в процессе межкультурного взаимодействия.
Компетенция была сформирована в процессе:
\begin{enumerate}
    \item изучения научных публикаций в международных изданиях и подходов к моделированию гидродинамики, 
        применяемых в научно-исследовательских организациях разных стран;
    \item участия в региональных научных конференциях;
    \item изучения публичных страниц ведущих исследователей в области моделирования движения жидкости методом SPH.
\end{enumerate}

УК-6 Способен определять и реализовывать приоритеты собственной деятельности и способы ее совершенствования на основе самооценки.
Компетенция была сформирована в процессе:
\begin{enumerate}
    \item планирования этапов разработки программного обеспечения для моделирования методом SPH с учётом сроков и других задач;
    \item регулярного анализа промежуточных результатов и корректировки направлений исследования;
    \item оценки эффективности выбранных алгоритмов и библиотек, их замены при необходимости;
    \item критического анализа собственных публикаций и презентаций с целью улучшения подачи материала;
    \item постановки личных целей развития в области вычислительной гидродинамики и суперкомпьютерных технологий.
\end{enumerate}

ОПК-1 Способен решать актуальные задачи фундаментальной и прикладной математики.
Компетенция была сформирована в процессе:
\begin{enumerate}
    \item изучения и анализа математических основ метода SPH, включая уравнения Навье–Стокса;
    \item разработки и реализации численных алгоритмов для решения задач гидродинамики с использованием метода SPH;
    \item верификации и валидации численных моделей путём сравнения с аналитическими решениями и экспериментальными данными;
    \item решения прикладных задач моделирования динамики жидкости;
    \item публикации результатов исследований в научных статьях и докладах на конференциях, 
        включая обсуждение математических основ и точности предложенных методов.
\end{enumerate}

ОПК-2 Способен совершенствовать и реализовывать новые математические методы решения прикладных задач.
Компетенция была сформирована в процессе:
\begin{enumerate}
    \item анализа существующих математических методов моделирования гидродинамики;
    \item реализации и модификации метода SPH;
    \item реализации средств восстановления профиля жидкости и определения значения функции поля в точке;
    \item разработки и реализации параллельных алгоритмов для ускорения расчётов 
        на суперкомпьютерных системах с использованием технологий OpenMP и OpenCL.
\end{enumerate}

ОПК-3 Способен разрабатывать математические модели и проводить их анализ при решении задач в области профессиональной деятельности.
Компетенция была сформирована в процессе:
\begin{enumerate}
    \item разработки математической модели движения жидкости на основе метода SPH;
    \item анализа адекватности и устойчивости модели;
    \item применения разработанных численных моделей для исследования волн на поверхности жидкости;
    \item подготовки научных публикаций с описанием моделей и анализом их возможностей.
\end{enumerate}

ОПК-4 Способен комбинировать и адаптировать существующие информационно-коммуникационные технологии 
для решения задач в области профессиональной деятельности с учетом требований информационной безопасности.
Компетенция была сформирована в процессе:
\begin{enumerate}
    \item применения технологий параллельного программирования OpenMP и OpenCL;
    \item использования системы контроля версий Git;
    \item использования системы сборки CMake;
    \item использования IDE Microsoft Visual Studio 2022 и Visual Studio Code;
    \item использования компиляторов MSVC, GCC и Clang;
    \item использования языков Python и C++, их стандартных библиотек и некоторых популярных библиотек для данных языков программирования;
    \item разработки и применения средств экспорта результатов моделирования при помощи DualSPHysics;
    \item построения графиков с помощью сервиса Desmos;
    \item использования системы компьютерной вёрстки \LaTeX;
    \item администрирования компьютера с антивирусным программным обеспечением.
\end{enumerate}

ПК-1 Способен планировать и организовывать проведение научно-исследовательских и 
опытно-конструкторских разработок в области информационных технологий и информационных систем, 
исходя из анализа проблем, потребностей и возможностей.
Компетенция была сформирована в процессе:
\begin{enumerate}
    \item обсуждения работы с научным руководителем;
    \item определения требований к пакету программ для суперкомпьютерного моделирования динамики жидкости;
    \item планирования этапов разработки программного обеспечения для моделирования методом SPH с учётом сроков и других задач;
    \item проектирования и реализации пакета программ для моделирования динамики жидкости;
    \item тестирования численных моделей и анализа промежуточных результатов;
    \item составления отчётов и презентаций в рамках практик, выпускной квалификационной работы и конференций.
\end{enumerate}

ПК-2 Способен проводить совещания рабочих групп, семинары, демонстрации и презентации по вопросам анализа информационных систем, 
научно-исследовательских и опытно-конструкторских разработок.
Компетенция была сформирована в процессе:
\begin{enumerate}
    \item обсуждения работы с научным руководителем;
    \item подготовки отчёта и презентации по выпускной квалификационной работе;
    \item подготовки отчётов по практическим работам;
    \item подготовки тезисов, презентаций и выступлений на региональных научных конференциях;
    \item обсуждения и ответов на вопросы во время выступлений на научных конференциях.
\end{enumerate}

ПК-3 Способен применять актуальную нормативную документацию при проведении 
научно-исследовательских и опытно-конструкторских разработок.
Компетенция была сформирована в процессе:
\begin{enumerate}
    \item подготовки текста выпускной квалификационной работы с учётом требований ГОСТ;
    \item составления библиографического списка;
    \item составления технической документации к разработанному пакету программ;
    \item изучения документации языков программирования C++, OpenCL C, Python и их библиотек;
    \item изучения документации для инструментов git, CMake;
    \item изучения документации для DualSPHysics;
    \item изучения документации \LaTeX и его пакетов. 
\end{enumerate}

ПК-4 Способен изучать новые предметные области, моделировать бизнес-процессы, 
разрабатывать информационные, объектные, документные модели производственных предприятий.
Компетенция была сформирована в процессе:
\begin{enumerate}
    \item составления блок-схем и диаграмм, описывающих проект;
    \item изучения научной литературы по теме исследования;
    \item изучения документации к используемым инструментам;
    \item применения системы контроля версий для сохранения, отслеживания и отката изменений в проекте;
    \item составления библиографического списка;
    \item проектирования и разработки ПО для моделирования динамики жидкости численным методом SPH;
    \item тестирования и анализа результатов моделирования для определения границ допустимых значений свободных параметров.
\end{enumerate}

ПК-5 Способен управлять проектами в области создания и применения информационных технологий и информационных систем, 
управлять спорами и конфликтами.
Компетенция была сформирована в процессе:
\begin{enumerate}
    \item обсуждения работы с научным руководителем;
    \item планирования этапов разработки программного обеспечения для моделирования методом SPH с учётом сроков и других задач;
    \item определения требований к пакету программ для моделирования динамики жидкости;
    \item проектирования, разработки и тестирования ПО для моделирования динамики жидкости методом SPH;
    \item использования сторонних средств моделирования и инструментов разработки.
\end{enumerate}
